\section{Lean target signatures}\label{app:signatures}

This appendix records the intended declarations independently of implementation details.  Notation is
adapted to the repository; exact elaboration may require minor implicit-argument adjustments.

\subsection{Plan kernel}

\begin{lstlisting}[language=Lean]
namespace LeanProofs.IPRA.TAPIR

open LeanProofs.ConciseIPRA
open Syn
open scoped Fm
open ConciseIPRA.PTm (zero succ)

inductive Plan : {n : Nat} -> List (Fm PTm n) -> Fm PTm n -> Type
  | prf {n} {Gamma : List (Fm PTm n)} {phi : Fm PTm n}
      (h : Gamma |- [PRAAx] phi) : Plan Gamma phi
  | weaken {n} {Gamma Delta : List (Fm PTm n)} {phi : Fm PTm n}
      (p : Plan Gamma phi) (s : Gamma subset Delta) : Plan Delta phi
  | subst {n m} {Gamma : List (Fm PTm n)} {phi : Fm PTm n}
      (p : Plan Gamma phi) (sigma : Fin n -> PTm m) :
      Plan (Gamma.map (fun psi => psi.subst sigma)) (phi.subst sigma)
  | cut {n} {Gamma Delta : List (Fm PTm n)} {phi : Fm PTm n}
      (p : Plan Gamma phi)
      (s : forall psi, psi in Gamma -> Plan Delta psi) : Plan Delta phi
  | botE ...
  | andI ... | andE1 ... | andE2 ...
  | orI1 ... | orI2 ... | orE ...
  | impI ... | impE ...
  | allI ... | allE ... | exI ... | exE ...
  | eqRefl ... | eqSubst ...
  | qfInd {n} {Gamma : List (Fm PTm n)}
      (theta : Fm PTm (n+1)) (hq : theta.QF)
      (p0 : Plan Gamma (theta.subst (Syn.inst zero)))
      (pS : Plan (theta :: Gamma.map Fm.shift)
        (theta.subst (Syn.repl (succ (var 0))))) :
      Plan Gamma (Fm.all theta)

def Plan.run {n} {Gamma : List (Fm PTm n)} {phi : Fm PTm n} :
    Plan Gamma phi -> (Gamma |- [PRAAx] phi)
\end{lstlisting}

The supplied candidate includes all displayed logical cases.  A smaller first merge could omit the
logical wrappers and retain only \code{prf}, structural operations, \code{exE}, and \code{qfInd}; the
schedule theorem needs only \code{prf} and \code{cut}.

\subsection{Schedule}

\begin{lstlisting}[language=Lean]
inductive Schedule {n : Nat} (R : List (Fm PTm n)) :
    List (Fm PTm n) -> Type
  | nil : Schedule R []
  | cons {outs : List (Fm PTm n)} {phi : Fm PTm n}
      (tail : Schedule R outs)
      (local : Plan (outs ++ R) phi) : Schedule R (phi :: outs)

def Schedule.proveAll
    (s : Schedule R outs) : forall psi, psi in outs -> Plan R psi

def Schedule.head
    (s : Schedule R (phi :: outs)) : Plan R phi :=
  s.proveAll phi (by simp)

theorem Schedule.sound
    (s : Schedule R (phi :: outs)) : R |- [PRAAx] phi :=
  Plan.run s.head
\end{lstlisting}

\subsection{Certificate macros}

\begin{lstlisting}[language=Lean]
def Plan.seal
    (hex : Plan Gamma (Fm.ex theta))
    (body : Plan (theta :: Gamma.map Fm.shift) psi.shift) :
    Plan Gamma psi := Plan.exE hex body

def Plan.closeResiduals
    (p : Plan R phi)
    (s : forall psi, psi in R -> Plan Delta psi) : Plan Delta phi :=
  Plan.cut p s

theorem boundedName
    (f : PTm 2)
    (h : Gamma |- [PRAAx]
      Fm.ex (lt (var 0) (Syn.shift N) /\
        neg (inst2 f (var 0) (Syn.shift x) = zero))) :
    Gamma |- [PRAAx]
      lt (wit f N x) N /\
      neg (inst2 f (wit f N x) x = zero) :=
  BoundedReflection.bounded_choice f h
\end{lstlisting}

\subsection{Generic-arity search}

\begin{lstlisting}[language=Lean]
namespace SearchK

variable {k n : Nat}

/-- Body variables: 0 = candidate, 1..k = parameters. -/
def liftBody (f : PTm (k+1)) : PTm (k+2) :=
  Syn.subst (Fin.cases (var 0) (fun j => var j.succ.succ)) f

/-- Step variables: 0 = N, 1 = old least zero, 2.. = parameters. -/
def muStep (f : PTm (k+1)) : PTm (k+2) :=
  var 1 + eqb (var 1) (var 0) * sg (liftBody f)

def muF (f : PTm (k+1)) : PR (k+1) := recDef zero (muStep f)

def mu (f : PTm (k+1)) (N : PTm n)
    (xs : Fin k -> PTm n) : PTm n :=
  app (muF f) (Fin.cases N xs)

theorem mu_zero ...
theorem mu_succ ...
theorem search_spec ...
theorem search_found ...
theorem search_minimal ...

end SearchK
\end{lstlisting}

The exact type of \code{zero} in \code{muF} is \code{PTm k}; Lean may need an annotation.  The
renaming shown is schematic; because of \code{Fin} coercions, a helper proved by \code{Fin.cases} is
recommended.

\subsection{Reflected formulas}

\begin{lstlisting}[language=Lean]
structure Reflected {n : Nat} (phi : Fm PTm n) where
  bit : PTm n
  cases : ([] : List (Fm PTm n)) |- [PRAAx]
    ((bit = zero) \/ (bit = one))
  sound : ([] : List (Fm PTm n)) |- [PRAAx]
    ((bit = one) --> phi)
  complete : ([] : List (Fm PTm n)) |- [PRAAx]
    (phi --> (bit = one))

namespace Reflected

def eq (a b : PTm n) : Reflected (a = b) := ...
def nonzero (a : PTm n) : Reflected (neg (a = zero)) := ...
def and (A : Reflected phi) (B : Reflected psi) :
    Reflected (phi /\ psi) := ...
def or (A : Reflected phi) (B : Reflected psi) :
    Reflected (phi \/ psi) := ...
def neg (A : Reflected phi) : Reflected (neg phi) := ...
def imp (A : Reflected phi) (B : Reflected psi) :
    Reflected (phi --> psi) := ...

def ball (body : Reflected theta) (N : PTm n) :
    Reflected (Fm.all (lt (var 0) (Syn.shift N) --> theta)) := ...

def bex (body : Reflected theta) (N : PTm n) :
    Reflected (Fm.ex (lt (var 0) (Syn.shift N) /\ theta)) := ...

end Reflected
\end{lstlisting}

The first version of \code{ball}/\code{bex} should be restricted to a body whose bit is explicitly a
term in candidate plus parameters.  A later family record can package the binder arithmetic more
elegantly.

\subsection{Conditional complete-list theorem}

\begin{lstlisting}[language=Lean]
structure BoundedPresentation (k : Nat) where
  sol : PTm (k+1)       -- variables: code, parameters
  height : PTm 1
  codeBound : PTm (k+1) -- H, parameters -> C
  solReflected : ...

theorem completeScan
    (P : BoundedPresentation k)
    (hBound : Gamma |- [PRAAx]
      Fm.all (Sol P --> le (height P) (Syn.shift H)))
    (hCode : Gamma |- [PRAAx]
      Fm.all (le (height P) (Syn.shift H) -->
        lt (var 0) (codeBound P H xs))) :
    Gamma |- [PRAAx] CompleteList P H xs := ...
\end{lstlisting}

The report recommends beginning with characteristic membership rather than a packed list code.  A
list encoder can be layered on after soundness and completeness of the bounded characteristic
function compile.
